%! AP Statistics — Unit 4, Lesson 4.10 — Student Packet Cover Sheet
\documentclass[10pt]{article}
\usepackage{apstats-article}
\usepackage{apstats-boxes}
\usepackage{ltablex}
\keepXColumns

\begin{document}

% Full-bleed navy banner
\begin{tikzpicture}[remember picture, overlay]
  \fill[navy] (current page.north west)
    rectangle ([yshift=-0.9in]current page.north east);
\end{tikzpicture}
\vspace{-0.8in}

\begin{center}
  {\color{white}\textbf{\LARGE AP Statistics}}\\[4pt]
  {\color{skymid}\large\textbf{Unit 4: Probability, Random Variables, and Probability Distributions}}\\[6pt]
  {\color{white}\textbf{\large Lesson 4.10 \quad Introduction to the Binomial Distribution}}
\end{center}

\vspace{0.22in}
\namedateperiod
\vspace{0.18in}

\begin{learningtargetbox}
\small
\begin{enumerate}[leftmargin=1.5em, itemsep=5pt]
  \item \ldots{} determine whether a random variable is binomial by checking the four
        BINS conditions (Binary, Independent, fixed $N$, Same $p$) and stating
        $X \sim B(n, p)$ when all conditions are met.
  \item \ldots{} compute a single binomial probability $P(X = k)$ using the formula
        $\binom{n}{k}p^k(1-p)^{n-k}$ and verify with \texttt{binompdf}$(n,p,k)$.
\end{enumerate}
\end{learningtargetbox}

\vspace{0.14in}

\begin{tocbox}
\small
\renewcommand{\arraystretch}{1.5}
\begin{tabularx}{\linewidth}{@{} c l X r @{}}
\rowcolor{sky}
\textbf{\#} & \textbf{Component} & \textbf{Description} & \textbf{Score} \\
\midrule
1 & Warm-Up        & Spiral review: discrete RVs, combinations, probability rules & \blank{1.2cm} \\
2 & Guided Notes   & BINS conditions; binomial formula; \texttt{binompdf}          & \blank{1.2cm} \\
3 & Group Activity & Identifying binomial settings; computing $P(X=k)$             & \blank{1.2cm} \\
4 & Exit Ticket    & Light-bulb defect scenario                                    & \blank{1.2cm} \\
5 & Homework       & Four binomial scenarios with formula setup                    & \blank{1.2cm} \\
\midrule
  & \multicolumn{2}{r}{\textbf{Total}} & \blank{1.2cm} \\
\end{tabularx}
\end{tocbox}

\vspace{0.14in}

\begin{remindbox}
\small
The \textbf{most important idea in Lesson 4.10}: before using any binomial formula,
check \textbf{all four BINS conditions} in writing.

\vspace{6pt}
\begin{tabularx}{\linewidth}{@{} >{\centering\arraybackslash\columncolor{goldbg}}p{0.06\linewidth}
                                    X @{}}
\textbf{\large B} &
\textbf{Binary outcomes.}\enspace
Every trial must produce exactly one of two outcomes (success or failure).
Clearly define which outcome is ``success.'' \\[6pt]
\textbf{\large I} &
\textbf{Independent trials.}\enspace
The outcome of one trial does not affect others.
When sampling without replacement, use the $10\%$ condition:
$n \le 0.10 \times N$. \\[6pt]
\textbf{\large N} &
\textbf{Fixed number of trials.}\enspace
The total number of trials $n$ is set in advance --- you do not
keep going until a success occurs. \\[6pt]
\textbf{\large S} &
\textbf{Same probability.}\enspace
The probability of success $p$ is the same on every trial.
\end{tabularx}
\end{remindbox}

\end{document}
